\section{Conclusion} \label{sec:56-research03-conclusion}

\begin{foldable} % Conclusion
  We propose TAKE, the first text \ac{DD} method to simultaneously satisfy all three desiderata: \textit{knowledge-based weighting} via trajectory-integrated influence scores that correct the \acl{HSB} and direct the budget toward informative samples; a \textit{dataset-level distillation objective} via discrete Optimal Transport that globally aligns the selected corpus with the source distribution; and \textit{human-readable output} via selection from \ac{LLM}-generated candidates, producing a distilled corpus that is inspectable, auditable, and transferable across architectures.
  These design choices address gaps that prior text \ac{DD} methods leave open: TAKE matches or exceeds baselines across six language benchmarks at extreme compression, with end-to-end theoretical formulation and guarantees.
  These results suggest that text dataset distillation can help address systemic cost pressures: reducing corpus size cuts cumulative training expenditure, lowers carbon footprint, and respects data-governance constraints without sacrificing transferability.
\end{foldable}
